% Conclusion
\section{Conclusion}
\label{sec:conclusion}

Is a ``truthful'' LLM truthful in all the same ways? Our analysis of 50 diverse language models provides a clear answer: No.

We presented the first systematic cross-benchmark correlation study for truthfulness evaluation, demonstrating that TruthfulQA, HaluEval, and FactScore measure three partially independent dimensions. Inter-benchmark correlations (r=0.42--0.58) are moderate---above unrelated-benchmark baselines but well below unity---indicating related but non-identical constructs. PCA confirms this structure: three components are required for 80\% variance, with no single factor dominating.

Our mechanism analysis revealed that:
\begin{itemize}
\item \textbf{TruthfulQA} measures misconception resistance, distinct from general knowledge (r=0.19 with MMLU vs.\ r=0.78 within MMLU)
\item \textbf{HaluEval} measures generation coherence, nearly orthogonal to misconception resistance (r=0.16 with TruthfulQA)
\item \textbf{FactScore} measures atomic factual precision, independent of both (r $\approx$ 0 with each)
\end{itemize}

The practical implication is direct: single-benchmark evaluation is insufficient. The divergent profile models---high MMLU but low TruthfulQA---demonstrate that knowing facts does not guarantee avoiding falsehoods. Comprehensive truthfulness assessment requires evaluating all three dimensions.

Looking forward, this work opens several directions. Cross-format validation could disentangle task paradigm effects from underlying construct differences. Longitudinal analysis could track whether multi-dimensional structure persists across model generations. Most importantly, understanding truthfulness as multi-dimensional enables targeted interventions: RLHF for misconception resistance, consistency training for generation coherence, retrieval augmentation for factual precision.

Truthfulness in LLMs is not one skill but several. Evaluation and improvement must treat it as such.
